\documentclass[11pt]{article}
% ======================================================================
%  examstyle.tex — EPFL-style exam layout for the weekly Time Series exams
%  Shared by every Exam_Week_NN(.tex / _Solutions.tex). Compiles with pdflatex.
% ======================================================================
\usepackage[utf8]{inputenc}
\usepackage[T1]{fontenc}
\usepackage[english]{babel}
\usepackage[a4paper,margin=2.4cm]{geometry}
\usepackage{amsmath,amssymb,amsthm,mathtools}
\usepackage{bm}
\usepackage{enumitem}
\usepackage{array}

\setlength{\parindent}{0pt}
\setlength{\parskip}{0.5em}
\emergencystretch=3em
\allowdisplaybreaks[1]

% ---------- math macros (same names as the rest of the project) ----------
\newcommand{\E}{\mathbb{E}}
\DeclareMathOperator{\Var}{Var}
\DeclareMathOperator{\Cov}{Cov}
\DeclareMathOperator{\Corr}{Corr}
\DeclareMathOperator*{\argmin}{arg\,min}
\DeclareMathOperator{\MSE}{MSE}
\DeclareMathOperator{\trace}{tr}
\newcommand{\Reals}{\mathbb{R}}
\newcommand{\Ints}{\mathbb{Z}}
\newcommand{\Comp}{\mathbb{C}}
\newcommand{\Nat}{\mathbb{N}}
\newcommand{\Prob}{\mathbb{P}}
\newcommand{\Tr}{^{\mathsf{T}}}
\newcommand{\conj}[1]{\overline{#1}}
\newcommand{\abs}[1]{\left\lvert #1 \right\rvert}
\newcommand{\norm}[1]{\left\lVert #1 \right\rVert}
\newcommand{\ind}{\mathbf{1}}
\newcommand{\eps}{\varepsilon}
\newcommand{\diff}{\nabla}
\newcommand{\acvs}{\gamma}
\newcommand{\acf}{\rho}
\newcommand{\pacf}{\alpha}
\newcommand{\sdf}{\mathcal{S}}
\newcommand{\Bsh}{B}
\newcommand{\iid}{\overset{\text{iid}}{\sim}}

\setlist[enumerate]{leftmargin=2em,topsep=2pt,itemsep=4pt}
\setlist[itemize]{leftmargin=1.6em,topsep=2pt,itemsep=2pt}

\newcounter{exo}

% ---------- exam cover header :  \examheader{exam title}{duration} ----------
\newcommand{\examheader}[2]{%
  \thispagestyle{empty}
  \begin{center}
    {\Large\bfseries Time Series}\\[3pt]
    Spring Semester 2026, EPFL\\
    \'Ecole Polytechnique F\'ed\'erale de Lausanne\\
    Prof.\ Dr.\ Sofia C.\ Olhede\\[7pt]
    \hrule height 0.8pt
    \vspace{9pt}
    {\large\bfseries Practice Exam --- #1}\\[4pt]
    {\normalsize Duration: #2 \quad$\cdot$\quad No calculator \quad$\cdot$\quad Answer on paper}\\
    \vspace{8pt}
    \hrule height 0.8pt
  \end{center}
  \vspace{14pt}
  \begin{tabular}{@{}p{8.2cm}@{}p{8cm}@{}}
    Name:\enspace\hrulefill & Surname:\enspace\hrulefill
  \end{tabular}\\[14pt]
  SCIPER (Student Card Number):\enspace\hrulefill
  \vspace{16pt}
}

% ---------- points table (fixed: 4 problems) ----------
\newcommand{\pointstable}{%
  \begin{center}
  \renewcommand{\arraystretch}{1.5}
  \begin{tabular}{|p{5cm}|p{4cm}|}
    \hline
    \textbf{Exercise} & \textbf{Points}\\\hline
    1 & \\\hline
    2 & \\\hline
    3 & \\\hline
    4 & \\\hline
    \textbf{Total} & \\\hline
  \end{tabular}
  \end{center}
}

% ---------- problem heading (exam: one problem per page) ----------
\newcommand{\problem}[1]{%
  \clearpage
  \stepcounter{exo}%
  \begin{center}\textbf{\large Problem \theexo\ (#1 points)}\end{center}
  \vspace{4pt}
}
% ---------- answer space: fill the statement page + one blank answer page ----------
% Put \answerpage at the END of each problem so every exercise gets its own page(s).
\newcommand{\answerpage}{\par\vfill\newpage\null\newpage}

% ---------- solutions document ----------
\newcommand{\solheader}[1]{%
  \begin{center}
    {\Large\bfseries Time Series --- Solutions}\\[3pt]
    {\large Practice Exam --- #1}\\[2pt]
    EPFL, MATH-342 \quad$\cdot$\quad Prof.\ S.\ C.\ Olhede\\[5pt]
    \hrule height 0.8pt
  \end{center}
  \vspace{10pt}
}
\newcommand{\solproblem}[1]{%
  \bigskip
  \stepcounter{exo}%
  {\large\bfseries Problem \theexo\ (#1 points)}\par\nobreak\vspace{2pt}
}
% Rappel de l'énoncé avant la solution (dans les corrigés)
\newcommand{\statementstart}{\par\vspace{1pt}{\bfseries\itshape Statement.}\par\nobreak\begingroup\small\itshape}
\newcommand{\statementend}{\par\endgroup\vspace{5pt}\hrule\vspace{6pt}{\bfseries Solution.}\par\nobreak\vspace{2pt}}

\begin{document}

\examheader{Week 7: Spectral Density \& Spectral Representation}{60 minutes}
\pointstable
\begin{center}\itshape Justify every answer. Throughout, $\eps_t$ is a mean-zero white noise of variance $\sigma^2$.\end{center}

% ---------------------------------------------------------------
\problem{5}
For a discrete-time stationary process $\{X_t\}_{t\in\Ints}$ with summable ACVS $\{\acvs_\tau\}\in\ell^1$ (sampling interval $\Delta=1$), the spectral density function (SDF) is $\sdf(f)=\sum_{\tau\in\Ints}\acvs_\tau\,e^{-2\pi i f\tau}$, with inverse $\acvs_\tau=\int_{-1/2}^{1/2}\sdf(f)\,e^{2\pi i f\tau}\,df$.

\textbf{(a)} Prove from the definition that any SDF is real-valued and even, $\sdf(-f)=\sdf(f)$, and write it in the cosine form $\sdf(f)=\acvs_0+2\sum_{\tau\ge1}\acvs_\tau\cos(2\pi f\tau)$. \hfill(1.5 pts)

\textbf{(b)} Show that $\acvs_0=\Var(X_t)=\int_{-1/2}^{1/2}\sdf(f)\,df$, and use this to interpret $\sdf(f)$ as a distribution of power across frequency. \hfill(1.5 pts)

\textbf{(c)} Take $\{X_t\}=\{\eps_t\}$ (white noise of variance $\sigma^2$). Compute $\sdf(f)$ directly from its ACVS and explain in one sentence why the result earns the name ``white''. \hfill(2 pts)
\answerpage

% ---------------------------------------------------------------
\problem{5}
This problem tests whether a candidate sequence is a valid autocovariance.

\textbf{(a)} \emph{(The validity criterion.)} State the criterion that decides whether a summable symmetric sequence $\{\acvs_\tau\}$ is the ACVS of \emph{some} stationary process, expressed through its Fourier transform $\sdf$. Why is non-negativity of $\sdf$ the right test? \hfill(1 pt)

\textbf{(b)} \emph{(A truncated ACVS.)} Let $\acvs_\tau=1$ for $\abs\tau\le K$ and $\acvs_\tau=0$ for $\abs\tau>K$, where $K$ is a positive integer. Compute $\sdf(f)$ in closed form (a Dirichlet kernel), then evaluate it at the Nyquist endpoint $f=\tfrac12$ to decide for which $K$ the sequence \emph{fails} to be a valid ACVS. \hfill(2 pts)

\textbf{(c)} \emph{(MA(1)-shape ACVS.)} Let $\acvs_0=1$, $\acvs_{\pm1}=\rho$, and $\acvs_\tau=0$ for $\abs\tau\ge2$, with $\rho\in\Reals$. Compute $\sdf(f)$ and determine the exact set of $\rho$ for which $\{\acvs_\tau\}$ is a valid ACVS. Relate the answer to the largest possible lag-$1$ autocorrelation of an MA(1). \hfill(2 pts)
\answerpage

% ---------------------------------------------------------------
\problem{5}
Let $\{X_t\}$ be the MA($q$) process $X_t=\sum_{j=0}^{q}\theta_j\eps_{t-j}$ with $\theta_0=1$ and $\eps_t$ white noise of variance $\sigma^2$.

\textbf{(a)} Recall (or derive) the ACVS $\acvs_\tau=\sigma^2\sum_{j=0}^{q}\theta_j\theta_{j+\abs\tau}$ for $\abs\tau\le q$ (and $0$ otherwise), writing it in the double-sum form $\acvs_\tau=\sigma^2\sum_{j=0}^{q}\sum_{k=0}^{q}\theta_j\theta_k\,\ind_{\{\tau=k-j\}}$. \hfill(1.5 pts)

\textbf{(b)} Starting from this double sum, derive the factored SDF
\[
\sdf(f)=\sigma^2\Bigl|\,\textstyle\sum_{j=0}^{q}\theta_j e^{-2\pi i j f}\,\Bigr|^2,
\]
and explain in one sentence why this form makes $\sdf(f)\ge0$ automatic. \hfill(2 pts)

\textbf{(c)} Specialise to the MA(1) $X_t=\eps_t-\theta\eps_{t-1}$. Give $\sdf(f)=\sigma^2\bigl(1+\theta^2-2\theta\cos(2\pi f)\bigr)$, and for $\theta>0$ locate the frequency of the maximum and of the minimum of $\sdf$, giving the two extreme values. \hfill(1.5 pts)
\answerpage

% ---------------------------------------------------------------
\problem{5}
\textbf{(Revision of Week 6: Fourier theory --- convolution \& aliasing.)} Work with the continuous-time Fourier transform $G(f)=\int_{-\infty}^{\infty}g(t)e^{-2\pi i f t}\,dt$ and convolution $(g*h)(t)=\int_{-\infty}^{\infty}g(t-y)h(y)\,dy$ ($g,h\in L^1$).

\textbf{(a)} \emph{(Convolution theorem.)} Prove that if $u=g*h$ then $U(f)=G(f)\,H(f)$, justifying carefully the interchange of integration (state the Fubini hypothesis you use). \hfill(2 pts)

\textbf{(b)} \emph{(Aliasing.)} Let $g\in L^1$ be continuous with transform $\mathcal G$, and let $G$ denote the transform of the sampled sequence $\{g(t)\}_{t\in\Delta\Ints}$. Show the aliasing identity $G(f)=\sum_{k\in\Ints}\mathcal G(f+k/\Delta)$, and define the Nyquist frequency. \hfill(2 pts)

\textbf{(c)} \emph{(A concrete alias.)} A pure cosine of frequency $f_0=0.8$ Hz is sampled at rate $1/\Delta=1$ Hz (so $\Delta=1$ s). To which frequency in the principal Nyquist band $[-\tfrac12,\tfrac12]$ does $f_0$ fold, and what is the lowest sampling rate that would resolve $f_0$ without aliasing? \hfill(1 pt)
\answerpage

\end{document}
